% !TeX program = pdflatex
\documentclass[11pt,a4paper]{article}

% Zentrale Einstellungen und Pakete
% ------------------------------------------------------------
% Sprache, Zeichencodierung und Anführungen
% ------------------------------------------------------------
\usepackage[ngerman]{babel}
\usepackage[utf8]{inputenc}
\usepackage[T1]{fontenc}
\usepackage{csquotes}

% ------------------------------------------------------------
% Seitenlayout: DIN A4, 2 cm Ränder
% ------------------------------------------------------------
\usepackage{geometry}
\geometry{
  top=2cm,
  bottom=2cm,
  left=2cm,
  right=2cm
}

% ------------------------------------------------------------
% Zeilenabstand und Feintypografie
% ------------------------------------------------------------
\usepackage{setspace}
\onehalfspacing
\usepackage{microtype}

% ------------------------------------------------------------
% Schrift: TeX Gyre Heros als Arial-ähnliche Schrift
% ------------------------------------------------------------
\usepackage{tgheros}
\renewcommand{\familydefault}{\sfdefault}

% ------------------------------------------------------------
% Überschriftenformatierung
% ------------------------------------------------------------
\usepackage{titlesec}

\titleformat{\section}
  {\bfseries\fontsize{12pt}{14pt}\selectfont}
  {\thesection}
  {1em}
  {}

\titleformat{\subsection}
  {\bfseries\fontsize{12pt}{14pt}\selectfont}
  {\thesubsection}
  {1em}
  {}

\titleformat{\subsubsection}
  {\bfseries\fontsize{12pt}{14pt}\selectfont}
  {\thesubsubsection}
  {1em}
  {}

\setcounter{secnumdepth}{3}
\setcounter{tocdepth}{3}

% ------------------------------------------------------------
% Absätze und Fußnoten
% ------------------------------------------------------------
\setlength{\parskip}{6pt}
\setlength{\parindent}{0pt}

\makeatletter
\renewcommand\footnotesize{\@setfontsize\footnotesize{10pt}{12pt}}
\makeatother

% ------------------------------------------------------------
% Seitenzahlen: zentriert unten, keine Kopfzeile
% ------------------------------------------------------------
\usepackage{fancyhdr}
\pagestyle{fancy}
\fancyhf{}
\fancyfoot[C]{\thepage}
\renewcommand{\headrulewidth}{0pt}

% ------------------------------------------------------------
% Grafiken, Tabellen und codebasierte Abbildungen
% ------------------------------------------------------------
\usepackage{graphicx}
\usepackage{xcolor}
\usepackage{tikz}
\usetikzlibrary{
  arrows.meta,
  backgrounds,
  calc,
  fit,
  matrix,
  positioning,
  shapes.geometric
}

% Einheitliche Farben der Fallstudiengrafiken
\definecolor{caseBlue}{HTML}{174A7E}
\definecolor{caseGreen}{HTML}{5A8F3D}
\definecolor{caseLight}{HTML}{F4F8FB}
\definecolor{caseLine}{HTML}{A9B4BF}

\tikzset{
  casePrimary/.style={
    draw=caseBlue,
    fill=caseBlue,
    text=white,
    rounded corners=3pt,
    align=center,
    font=\sffamily\bfseries\scriptsize,
    minimum height=8mm,
    inner sep=4pt
  },
  caseBox/.style={
    draw=caseBlue,
    fill=white,
    text=caseBlue,
    rounded corners=3pt,
    align=center,
    font=\sffamily\scriptsize,
    minimum height=8mm,
    inner sep=4pt
  },
  caseOptional/.style={
    draw=caseGreen,
    fill=caseGreen!10,
    text=caseGreen!55!black,
    rounded corners=3pt,
    align=center,
    font=\sffamily\scriptsize,
    minimum height=8mm,
    inner sep=4pt
  },
  caseGroup/.style={
    draw=caseLine,
    fill=caseLight,
    rounded corners=5pt,
    inner sep=6pt
  },
  caseFlow/.style={
    -{Latex[length=2.6mm,width=1.8mm]},
    draw=caseBlue,
    line width=0.9pt
  },
  caseSecondaryFlow/.style={
    -{Latex[length=2.4mm,width=1.7mm]},
    draw=caseLine,
    line width=0.8pt,
    dashed
  }
}

\usepackage{booktabs}
\usepackage{array}
\usepackage{tabularx}
\usepackage{float}
\usepackage{threeparttable}

% ------------------------------------------------------------
% Hyperlinks
% ------------------------------------------------------------
\usepackage[hidelinks]{hyperref}
\usepackage{url}

% ------------------------------------------------------------
% Literaturverwaltung mit biblatex-apa
% ------------------------------------------------------------
\usepackage[
  style=apa,
  backend=biber,
  sorting=nyt,
  giveninits=true,
  maxcitenames=2,
  maxbibnames=20,
  doi=true,
  url=true,
  isbn=true
]{biblatex}

\DeclareLanguageMapping{ngerman}{ngerman-apa}

\DefineBibliographyStrings{ngerman}{
  nodate = {o\adddot\addspace D\adddot}
}

\addbibresource{references.bib}


% ------------------------------------------------------------
% Metadaten
% Die Werte in eckigen Klammern vor der Abgabe ersetzen.
% Keine persönlichen oder institutionellen Angaben erfinden.
% ------------------------------------------------------------
\newcommand{\hochschule}{[Name der Hochschule]}
\newcommand{\arbeitstitel}{[Titel der Fallstudie zum Technologie-Template]}
\newcommand{\arbeituntertitel}{[Untertitel]}
\newcommand{\arbeitart}{[Art der Arbeit]}
\newcommand{\kurs}{[Kurs / Modul]}
\newcommand{\aufgabenstellung}{[Aufgabenstellung]}
\newcommand{\studiengang}{[Studiengang]}
\newcommand{\autorname}{[Vorname Nachname]}
\newcommand{\matrikel}{[Matrikelnummer]}
\newcommand{\tutor}{[Tutor / Betreuer]}
\newcommand{\abgabedatum}{[Monat Jahr]}

% Auf \literatureavailabletrue setzen, sobald echte Quellen zitiert werden.
\newif\ifliteratureavailable
\literatureavailablefalse

\begin{document}

% ------------------------------------------------------------
% Titelblatt
% ------------------------------------------------------------
\hypersetup{pageanchor=false}
\pagenumbering{gobble}
\thispagestyle{empty}

\begin{titlepage}
  \centering

  {\Large \hochschule \par}
  \vspace{1.2cm}

  {\Large \arbeitart \par}
  \vspace{1.0cm}

  {\bfseries\LARGE \arbeitstitel \par}
  \vspace{0.4cm}

  {\large \arbeituntertitel \par}
  \vspace{1.5cm}

  {\large \aufgabenstellung \par}
  \vspace{2.0cm}

  \begin{tabularx}{0.82\textwidth}{@{}lX@{}}
    \textbf{Kurs:}           & \kurs \\
    \textbf{Studiengang:}    & \studiengang \\
    \textbf{Autor:}          & \autorname \\
    \textbf{Matrikelnummer:} & \matrikel \\
    \textbf{Tutor:}          & \tutor \\
    \textbf{Abgabedatum:}    & \abgabedatum \\
  \end{tabularx}

  \vfill

  {\small
  Schriftliche Ausarbeitung im Rahmen der Fallstudie zum Kurs\\
  \kurs
  \par}
\end{titlepage}

\hypersetup{pageanchor=true}
\pagenumbering{Roman}
\setcounter{page}{1}

% ------------------------------------------------------------
% Verzeichnisse
% ------------------------------------------------------------
\clearpage
\tableofcontents

\clearpage
\renewcommand{\listfigurename}{Abbildungsverzeichnis}
\phantomsection
\addcontentsline{toc}{section}{Abbildungsverzeichnis}
\listoffigures

\clearpage
\renewcommand{\listtablename}{Tabellenverzeichnis}
\phantomsection
\addcontentsline{toc}{section}{Tabellenverzeichnis}
\listoftables

% ------------------------------------------------------------
% Abkürzungsverzeichnis
% ------------------------------------------------------------
\clearpage
\phantomsection
\section*{Abkürzungsverzeichnis}
\addcontentsline{toc}{section}{Abkürzungsverzeichnis}

\begin{tabularx}{\textwidth}{@{}lX@{}}
\textbf{Abk.} & \textbf{Bedeutung} \\
\hline
% Später verwendete Abkürzungen prüfen und nur tatsächlich
% verwendete Begriffe aufnehmen.
%
% Mögliche Kandidaten:
% API, ATP, CI, CLI, CSP, HTTP, JSON, ORM, SQL
\end{tabularx}

% ------------------------------------------------------------
% Hauptkapitel
% Nur diese Datei steuert Reihenfolge und Auswahl der Hauptkapitel.
% ------------------------------------------------------------
\clearpage
\pagenumbering{arabic}
\setcounter{page}{1}

\input{workspace/chapters/100_einleitung/100}
\input{workspace/chapters/200_anforderungen/200}
\input{workspace/chapters/300_architektur/300}
\input{workspace/chapters/400_projektprofile/400}
\input{workspace/chapters/500_tooling_workflow/500}
\input{workspace/chapters/600_qualitaet_verifikation/600}
\input{workspace/chapters/700_bewertung/700}
\input{workspace/chapters/800_einsatz_transfer/800}
\input{workspace/chapters/900_schluss/900}

% ------------------------------------------------------------
% Literaturverzeichnis
% Die Bibliographie bleibt bis zur Aufnahme echter Quellen leer.
% ------------------------------------------------------------
\clearpage
\ifliteratureavailable
  \printbibliography[heading=bibintoc,title={Literaturverzeichnis}]
\else
  \phantomsection
  \section*{Literaturverzeichnis}
  \addcontentsline{toc}{section}{Literaturverzeichnis}
\fi

% ------------------------------------------------------------
% Anhang
% Bei Bedarf später mit \appendix und app:-Labels ergänzen.
% ------------------------------------------------------------

\end{document}
